\section{Counting the exceptional characters and primes}
\label{sec:window-exceptions}

We prove the zero-density estimate and give the progression mean-square
reduction used in Section~\ref{sec:rh-common-windows}. The density
estimate counts primitive characters for which at least one zero
occurs. Each character contributes one to the count irrespective of
the multiplicity or number of such zeros.

\subsection{A detector for a single zero}

Put \(\sigma_0=399/400\). For \(Q\ge3\) consider primitive
characters \(\chi\) of conductor \(q\le Q\) with a zero \(\rho\)
satisfying \(\Re\rho>\sigma_0\), \(|\Im\rho|\le Q\).
There is an absolute constant \(C\) such that their number is at
most \(CQ^{1/4}\).

It suffices to treat large \(Q\), since a bounded conductor range
contains only finitely many characters. The character of conductor
one contributes at most one. For each of the others select one zero
in the box. These characters are nonprincipal, and the classical
nonvanishing theorem on \(\Re s\ge1\) gives \(\Re\rho<1\).
For \(Q\ge20\) set
\[
 X=\lceil Q\rceil,\quad Y=X^4,\quad
 J=8(\lfloor\log_2 X\rfloor+1),\quad N=2^J .
\]
Then \(X^8\le N\le256X^8\). Write \(P_\chi(s)=
\sum_{b\le N}\chi(b)b^{-s}\) and \(M_\chi(s)=
\sum_{a\le Y}\mu(a)\chi(a)a^{-s}\), with all sums starting at one.
Character orthogonality over a complete period bounds each prefix
sum of \(\chi\) by \(q\). In the half-plane \(\sigma=\Re s>0\)
the resulting Dirichlet series converges to \(L(s,\chi)\) by partial
summation and analytic continuation from \(\Re s>1\). More explicitly,
the tail is bounded by \(2(q+1)(1+|s|/\sigma)N^{-\sigma}\):
apply partial summation to \(t^{-s}\), use the prefix bound at the
endpoint and integrate \(|s|t^{-\sigma-1}\) to infinity.
At the selected zero this bounds \(|P_\chi(\rho)|\).
Now \(q+1\le2X\), \(1+|\rho|/\Re\rho\le3X\), and
\(N^{\Re\rho}\ge N^{15/16}\ge X^7\sqrt X\ge4X^7\).
As \(|M_\chi(\rho)|\le Y\), we have
\[
     |P_\chi(\rho)M_\chi(\rho)|
       \le\frac{2(2X)(3X)}{4X^7}X^4=\frac{3}{X}\le\tfrac12.
\]
The Möbius convolution makes this small product a detector. Define
\(a_n=\sum_{a\le Y,\,b\le N,\,ab=n}\mu(a)\).
Then \(a_1=1\), and for \(1<n\le Y\) the sum runs over every
divisor of \(n\), so \(a_n=0\). Consequently
\[
 P_\chi(\rho)M_\chi(\rho)=
     1+\sum_{Y<n\le YN}a_n\chi(n)n^{-\rho}.
\]
The last sum has magnitude at least \(1/2\). The \(J\) consecutive
blocks \((U,2U]\), \(U=Y2^j,\ 0\le j<J\), partition its range.
Select a block for which
\[
 \left|\sum_{U<n\le2U}a_n\chi(n)n^{-\rho}\right|
              \ge\frac{1}{2J}.
\]

\subsection{How many characters can use one block?}

The elementary divisor bound \(d(n)\le C_0 n^{1/400}\) applies
uniformly in \(Y,N\), since \(|a_n|\le d(n)\).
To obtain it, for sufficiently large primes \(p\) use
\(e+1\le2^e\le p^{e/400}\) for every \(e\ge0\).
For each of the finitely many smaller primes the supremum of
\((e+1)p^{-e/400}\) is finite. Their product gives \(C_0\).
Hence \(b_n=a_nn^{-\sigma_0}\) satisfies, for \(U\ge1\),
\[
 \sum_{U<n\le2U}|b_n|^2
   \le C_0^2 U^{1-2(199/200)}
   \le C_0^2 U^{-197/200}.
\]
Put \(C_1=C_0^2\). For the set \(T\) of characters assigned to this
block define \(v_\chi(n)=\chi(n)n^{-(\rho_\chi-\sigma_0)}\).
Each Gram matrix diagonal is at most \(U\). For distinct primitive
\(\chi,\chi'\) the character \(\chi\overline{\chi'}\) lifted to
modulus \(qq'\) is nonprincipal. Otherwise their inflations agree
on units and have the same primitive inducing character, forcing
the original characters, including conductors, to agree.
Its sum over a full period vanishes and every prefix has absolute
value at most \(qq'\). For \(w=\rho_\chi-\sigma_0+
\overline{\rho_{\chi'}-\sigma_0}\), partial summation gives
\[
 \left|\sum_{U<n\le2U}\chi(n)\overline{\chi'(n)}n^{-w}\right|
       \le 2qq'(1+|w|)\le8Q^3.
\]
Indeed \(\Re w>0\), \(|\Im w|\le2Q\), \(\Re w<2(1-\sigma_0)\);
on the interval the total variation of \(t^{-w}\) is at most
\(|w|\int_U^{2U}dt/t\le|w|\). Prefix differences cost at most
\(2qq'\). The sum of the absolute values in each Gram row is
therefore at most \(\Delta=U+8|T|Q^3\).

For clarity, put \(s_\chi=\sum_n b_n v_\chi(n)\),
\(Z=\sum_\chi|s_\chi|^2\), and \(w_n=\sum_\chi\overline{s_\chi}
v_\chi(n)\). Then \(Z=\sum_n b_n w_n\). By Cauchy--Schwarz,
\(Z^2\le\|b\|_2^2\sum_n|w_n|^2\). Expanding the second factor and
using \(2|s_\chi||s_{\chi'}|\le |s_\chi|^2+|s_{\chi'}|^2\)
bounds it by \(\Delta Z\), since the absolute Gram matrix is symmetric.
Division by \(Z\) when \(Z>0\), with the zero case immediate, gives the claimed bound.
The detector now gives
\[
 \frac{|T|}{4J^2}
 \le C_1 U^{-197/200}\{U+8|T|Q^3\}.
\]
Since \(U\ge X^4\ge Q^4\), the coefficient of \(|T|\) on the right
is at most \(8C_1Q^{-47/50}\), which is at most \(1/(8J^2)\)
for all sufficiently large \(Q\). Moving that term to the left yields
\[
                 |T|\le8C_1 J^2 U^{3/200}.
\]
There are \(J\) possible blocks, each with
\(U\le YN\le256X^{12}\le2^{20}Q^{12}\). For all sufficiently
large \(Q\), \(J\le32\log Q\) and \(J^3\le Q^{7/100}\).
Including the possible character of conductor one, the count is
therefore bounded by
\[
 1+8C_1 (2^{20})^{3/200}
           Q^{7/100+12(3/200)}\ll Q^{1/4}.
\]
The classical inputs in this argument are orthogonality of
Dirichlet characters, uniqueness of primitive induction, analytic
continuation and nonvanishing on the line \(\Re s=1\) for
nonprincipal Dirichlet characters. Each bound on the detector,
the block Gram matrix and the final count has been derived above.

In a prime stage with upper endpoint \(V\) and base conductor
at most \(F\le\sqrt V\), put \(Q=FV\). A candidate \(p>\sqrt V\)
failing the zero-free condition yields a primitive offending twist
of conductor \(fp\) for \(1\le f\le F\). If the twists for \(p,p'\)
coincide, then \(fp=f'p'\) and \(p>F\ge f'\) give \(p=p'\).
There are at most \(O(Q^{1/4})=O(V^{3/8})\) failures. Against the
source's separate prime-progression lower bound \(V^{9/10}\), this
is a fraction \(O(V^{-21/40})\).

\subsection{The quadratic mean square on the actual progressions}

For this part write \(Y\) for the lower endpoint denoted \(U\) in
Section~\ref{sec:rh-common-windows}, and \(V\) for the upper cutoff
on candidate primes. Choose \(0<\epsilon\le1\) and \(K\ge1\).
Set \(w=Y^{\epsilon/5}\), \(Z=(Y^K)^{400}\), \(N=\lfloor Z\rfloor\)
and \(L=\log Y\). The base progression has modulus \(m_0>0\)
divisible by \(8\), and a reduced class \(a_0\). Assume
\[
 \tfrac14 L^2\log m_0\le\log V,\qquad
 Y^{\epsilon/6}\le\log V\le Y^{\epsilon/5}L^2,\qquad
 p\mid m_0\ \Longrightarrow\ p\le w\ \text{or}\ p>Z
\]
for prime \(p\). Let \(\mathcal P=\{q:\sqrt V<q\le V,\,
q\text{ prime},\ q\equiv a_0\bmod m_0\}\).
For \(Y\) sufficiently large in terms of \(\epsilon,K\),
every \(q\in\mathcal P\) exceeds \(N^2\), since
\(1600K\log Y\le Y^{\epsilon/6}\le\log V\).
For \(Y\le y\le Y^K\) put
\[
 c(n)=\Lambda(n)\min(1,y/n),\qquad
 S_q=\sum_{1\le n\le N} c(n)\kappa(n)\left(\frac{n}{q}\right).
\]
We show that
\begin{equation}
 \sum_{q\in\mathcal P}|S_q|^2\le C y\log y\,|\mathcal P|
 \qquad (|\kappa(n)|\le1).
 \label{eq:window-progress-variance}
\end{equation}
The absolute constant \(C\) is independent of \(\epsilon,K\);
the lower threshold on \(Y\) may depend on them.

The input about primes in progressions is the following instance
of \texttt{twistK}. For all but at most one odd squarefree
\(s>1\) coprime to \(m_0\) and bounded by \(Y^{800K}\), one has
\begin{equation}
 \left|\sum_{q\in\mathcal P}\left(\frac{q}{s}\right)\right|
       \le\eta|\mathcal P|,\qquad \eta=Y^{-(K+1)} .
 \label{eq:window-jacobi}
\end{equation}
We give its classical inputs explicitly. The relative
prime number theorem in arithmetic progressions in the source
has main term \(\operatorname{li}(V)/\varphi(M)\), with a correction
\(-\chi_1(b)\operatorname{li}(V^{\beta_1})/\varphi(M)\)
when the real exceptional zero is present. Here
\(M=m_0s\), \(b\) is reduced, and the relative error is
\[
 O\!\left(e^{-c\log V/\log M}+e^{-\sqrt{c\log V}}\right)
 \quad\text{when}\quad C_0\log M\log(\log M+2)\le\log V .
\]
The zero-free region and the Landau--Page uniqueness bound are
applied across the moduli \(M\le H=m_0Y^{800K}\).
Summing the main term against the Jacobi character across the
classes \(b\equiv a_0\bmod m_0\) gives zero by orthogonality.
An exceptional term survives only if the exceptional character
agrees with the Jacobi character on units congruent to \(1\bmod m_0\).
Two such characters with \(1-\beta_1<c'/\log H\) induce the
same primitive character by Landau--Page. Hence the two Jacobi
characters coincide on those units. By the Chinese remainder
theorem a differing odd prime factor in their squarefree moduli
would allow a quadratic nonresidue there and ones at all other
factors and modulo \(m_0\), forcing opposite values.
There is thus at most one exceptional modulus \(s_*\).
When \(1-\beta_1\ge c'/\log H\) the surviving correction
relative to \(\operatorname{li}(V)\) is bounded by
\(4\log V\exp(-c'\log V/\log H)\). For fixed \(\epsilon,K\)
the hypotheses imply \(\log V/\log H\gg(\log Y)^2\) and
\(\log\log V=O(\log Y)\), while \(\sqrt{\log V}\) grows faster
than every fixed power of \(\log Y\). These errors are
\(o(Y^{-(K+1)})\). Removing primes up to \(\sqrt V\)
costs at most \(V^{1/2}\le V^{-2/5}|\mathcal P|\), by the
separate progression-count input \(|\mathcal P|\ge V^{9/10}\).
The relative-error estimate also applies to the unweighted count
of the classes, which is the denominator in \eqref{eq:window-jacobi}.
This is the role of the relative form of the prime number theorem.
The source proves the uniform comparison in \texttt{twistK}
and invokes the relative prime number theorem, the zero-free region
and the exceptional-zero bound there.

Expand the square of \(S_q\) and use \(|\kappa|\le1\).
For \(h=nn'\), define \(s(h)\) to be the product of primes
outside \(m_0\) having odd exponent in \(h\).
Quadratic reciprocity and its supplement for \(2\) give
\[
 \left|\sum_{q\in\mathcal P}\left(\frac{nn'}q\right)\right|
     =\left|\sum_{q\in\mathcal P}\left(\frac{q}{s(nn')}\right)\right|.
\]
Indeed every prime factor of \(nn'\) is less than \(q\);
square factors contribute one. The sign in reciprocity depends
only on \(q\bmod4\), the factor from \(2\) only on \(q\bmod8\),
and the factors from primes dividing \(m_0\) are fixed by the
progression. Their product is independent of \(q\).
The bound \(s(nn')\le N^2\le Y^{800K}\) permits
\eqref{eq:window-jacobi} at all pairs except those with
\(s(nn')=1\) or \(s(nn')=s_*\).

Write \(A_p=\sum_{p^j\le N,\,j\ge1}c(p^j)\),
\(B=\sup_p A_p\), and let \(R\) be the sum of \(c(p^j)\)
over powers for which \(j\) is even or \(p\mid m_0\).
Only prime powers have nonzero coefficient.
For \(s(nn')=1\), either both powers have the same base prime or
both have trivial squarefree part outside \(m_0\). Their total
weight is at most \(\sum_p A_p^2+R^2\).
For a specified \(s_*>1\), if it is prime, pairs with different
base primes have one factor with outside part \(s_*\) and the
other with trivial outside part, giving at most \(2BR\); pairs
with the same base prime have weight at most \(B^2\).
If \(s_*\) has two distinct prime factors the only possible
bases are those two, giving at most \(2B^2\).
Three or more prime factors give no pairs.
Use the sum \(2B^2+2BR\) in all cases. We have obtained
\begin{equation}
 \sum_{q\in\mathcal P}|S_q|^2\le|\mathcal P|
 \left\{\sum_p A_p^2+R^2+2B^2+2BR+
                      \eta\left(\sum_{n\le N}c(n)\right)^2\right\}.
 \label{eq:window-pair-split}
\end{equation}

Here are the size estimates that keep the constant absolute.
For \(y\ge3\), \(A_p\le3\log y\). If \(p\le y\) sum the terms
through \(\lfloor\log y/\log p\rfloor\) and then the decreasing
geometric tail, whose sum in units of \(\log p\) is at most \(2\).
If \(p\ge y\), use \(A_p\le 2y\log p/p\le2\log y\) by
monotonicity of \((\log t)/t\) for \(t\ge3\).
The prime diagonal satisfies
\(\sum_p A_p^2\le C_3 y\log y\). Indeed the primes up to
\(\sqrt y\) contribute at most \(9\sqrt y(\log y)^2\);
on \(\sqrt y<p\le y\) use
\(A_p^2\le18\log y\log p\) and
\(\sum_{p\le y}\log p\le y\log4\).
For \(p>y\) use \(A_p^2\le4y^2(\log p)^2/p^2\).
On \(t<p\le2t\), the sum of \((\log p)^2/p^2\) is at most
\(2(\log4)\log(2t)/t\), by the same Chebyshev estimate. Summing
the dyadic bounds \(t=2^j y\) gives \(O(\log y/y)\) and proves
the diagonal estimate. The Chebyshev bound follows, for example,
from the primorial bound \(\prod_{p\le n}p\le4^n\).

For \(R\), the total from even powers is at most
\(\sum_{a\le N}\Lambda(a)\min(1,y/a^2)\).
Use \(\sum_{a\le t}\Lambda(a)\ll t\), the elementary Chebyshev
bound, on \(a\le\sqrt y\) and on successive dyadic intervals above
\(\sqrt y\). The first sum is \(O(\sqrt y)\), and interval \(j\)
in the tail is \(O(2^{-j}\sqrt y)\).
All remaining base primes in \(R\) divide \(m_0\) and are at most
\(N\), hence at most \(w\). There are at most \(w\) of them,
each contributing at most \(3\log y\).
Thus \(R\le C_1(\sqrt y+w\log y)\).
Since \(w\le y^{1/5}\), the elementary inequality
\(\log y\le(10/3)y^{3/10}\) gives \(R\le5C_1\sqrt y\).
Finally, partial summation of the Chebyshev bound gives
\(\sum_{n\le N}c(n)\le y\sum_{n\le N}\Lambda(n)/n\le C_4 y\log Z\)
in the range at issue. For fixed \(K\) and all sufficiently large
\(Y\), \((\log Z)^2=(400K\log Y)^2\le Y\).
Then \(\eta y^2(\log Z)^2\le y\), since \(y\le Y^K\).
All five expressions in braces in \eqref{eq:window-pair-split}
are \(O(y\log y)\). More precisely with \(B\le C_2\log y\),
the source takes
\[
 C=C_3+25C_1^2+4C_2^2+20C_1C_2+C_4^2.
\]
This proves \eqref{eq:window-progress-variance}, including its
uniformity for arbitrary complex multipliers.

At the first selection take \(m_0=8,\ a_0=3,\ V=e^w\).
At the second take \(m_0=16q\prod_{2<p\le w}p\), the reduced
class of Section~\ref{sec:rh-common-windows}, and
\(V=e^{w(\log Y)^2}\). In this second stage \(q>\sqrt{e^w}>Z\)
for large \(Y\), so adjoining \(q\) to \(m_0\) preserves the
condition on prime divisors. The modulus bound follows from
\(\log q\le w\) and \(\sum_{p\le w}\log p\le w\log4\),
which give \(\log m_0\le4w\) for large \(Y\).
Thus the same \(K\) and lower endpoint \(Y\) apply in both stages.
The moving-point estimate and finite union over characters in
Section~\ref{sec:rh-common-windows} then give the required
exceptional fraction with \(\tau=\epsilon/2\) on all of
\([Y,Y^K]\).

\paragraph{Source correspondence.}
The density proof is
\texttt{CountZAll.zero\_density\_box}, with the exact detector in
\texttt{CountZDetect}, the coefficient square estimate in
\texttt{CountZMoll}, the interval character sums in
\texttt{CountZPair} and \texttt{CountZTwist}, the Gram bound in
\texttt{CountZBlock} and \texttt{CountZDual}, and the parameter
choice in \texttt{CountZNum} and \texttt{CountZFinal}. These are
the modules under
\srcid{Formalism.GL3MeanSquares.GenusMeanAsymptotics.JacobiCharacterSums}.
For the second estimate \srcid{QH.meanSqK} uses
\texttt{twistK}, \texttt{MeanSqAbs}, \texttt{SumMinLeDiv},
\texttt{SPsqLe}, \texttt{RRLe} and \texttt{SumCNLe} under
\srcid{Development.AutomorphicArithmetic.MeanSquareBounds}.
The relative prime number theorem is invoked as
\texttt{pnt\_progressions\_relative\_weak\_holds}, the zero-free
region as \texttt{zFRDirichlet}, the exceptional-zero estimate as
\texttt{excBound}, and the rational-prime count as
\texttt{count\_a}. The explicit formula that converts these
finite prime sums to centered sums of the twists is a separate
input in \texttt{BoxOfS}.
